% !PRINTED 638
\weberpage{638}
%<<CONT
contained function, when the $x_1, x_2, \ldots, x_n$ are integral rational
functions, an integral function in $\Omega$. But it is not thereby
said that conversely, in the form (2) with \gesperrt{integral} coefficients
 $x_i$, all integral functions in $\Omega$ are contained.

If, then, there are integral functions in the form (2) in which the
$x_i$ are not all integral rational functions, we can find an
integral function:

\[
\frac{x_1 \omega_1 + x_2 \omega_2 + \cdots + x_n \omega_n}{z - c},
\]

in which the $x_i$ are integral and not all divisible by $z - c$.
If we reduce the $x_i$ to their constant residues (with respect to $z - c$),
there results an integral function $\omega$ in the form:

\weq{3}{\omega = \frac{c_1 \omega_1 + c_2 \omega_2 + \cdots + c_n \omega_n}{z - c},}

in which the constants $c_1, c_2, \ldots, c_n$ certainly do not all vanish.
If, say, $c_1$ is different from zero, we can express $\omega_1$
by $\omega, \omega_2, \ldots, \omega_n$ and obtain a new integral
 basis of $\Omega$:

\weq{4}{\omega, \omega_2, \ldots, \omega_n.}

But for the discriminant there results, by § 173, (6), the
relation:

\weq{5}{\varDelta(\omega, \omega_2, \ldots, \omega_n) = \frac{c_1^2}{(z-c)^2} \varDelta(\omega_1, \omega_2, \ldots, \omega_n).}

Both discriminants are integral rational functions of $z$,
but the discriminant of the basis (4) is of \gesperrt{lower degree
than the discriminant of the basis} (1).

If we continue in this way to lower the degree of the
 discriminant, we must finally arrive at a basis $\omega_1, \omega_2, \ldots, \omega_n$
consisting of integral functions and having the property that in the form (2), with
integral rational $x_i$, all and only the integral functions in $\Omega$ are expressed.

8. \gesperrt{Definition: An integral basis $\omega_1, \omega_2, \ldots, \omega_n$
of the field $\Omega$ is called a minimal basis if in the
form}

\weq{6}{\omega = x_1 \omega_1 + x_2 \omega_2 + \cdots + x_n \omega_n}

\gesperrt{all integral functions of the field $\Omega$ can be represented
with integral rational $x_1, x_2, \ldots, x_n$.}

Such a minimal basis therefore always exists.
